% SPDX-License-Identifier: MIT
% Copyright (c) 2026 Evgeny Baulin
% Optimization Methods, Seminar 04: theory handout (English).
% Generated by checks/04. Accelerated Gradient Methods.ipynb from web/04/js/content-*.js and content/students.json. Do not edit by hand.
% Compile with: latexmk -xelatex -interaction=nonstopmode -halt-on-error Theory_en.tex
\documentclass[11pt,a4paper]{article}

\usepackage{fontspec}
\setmainfont{cmunrm.otf}[BoldFont=cmunbx.otf, ItalicFont=cmunti.otf, BoldItalicFont=cmunbi.otf]
\setsansfont{cmunss.otf}[BoldFont=cmunsx.otf, ItalicFont=cmunsi.otf, BoldItalicFont=cmunso.otf]
\setmonofont{cmuntt.otf}[BoldFont=cmuntb.otf, ItalicFont=cmunit.otf, BoldItalicFont=cmuntx.otf]
\usepackage{polyglossia}
\setdefaultlanguage{english}

\usepackage[a4paper,left=2.3cm,right=2.3cm,top=2.4cm,bottom=2.4cm,headheight=14pt]{geometry}
\usepackage{amsmath,amssymb}
\usepackage{graphicx}
\usepackage{array,booktabs,longtable}
\usepackage{enumitem}
\usepackage{xcolor}
\usepackage[most]{tcolorbox}
\usepackage{titlesec}
\usepackage{fancyhdr}
\usepackage{needspace}
\usepackage[expansion=false]{microtype}
\usepackage[colorlinks=true,linkcolor=hseblue,urlcolor=hseblue,citecolor=hseblue]{hyperref}

\definecolor{hseblue}{RGB}{16,45,105}
\definecolor{hseaccent}{RGB}{198,40,40}
\definecolor{hsegray}{RGB}{108,114,126}
\definecolor{hselight}{RGB}{240,243,248}
\definecolor{hsewarm}{RGB}{253,244,232}

\hypersetup{pdftitle={Optimization Methods. Seminar 04. Theory handout}, pdfauthor={Evgeny Baulin}}

\titleformat{\section}{\color{hseblue}\Large\bfseries}{\thesection}{0.6em}{}
\titleformat{\subsection}{\color{hseblue}\large\bfseries}{\thesubsection}{0.6em}{}
\titlespacing*{\section}{0pt}{16pt plus 4pt minus 2pt}{8pt plus 2pt}
\titlespacing*{\subsection}{0pt}{12pt plus 3pt minus 2pt}{6pt plus 2pt}
% a heading keeps a few lines under it on its page: a breakable box right after it would otherwise leave it alone
\pretocmd{\section}{\Needspace{7\baselineskip}}{}{\errmessage{cannot patch section}}
\pretocmd{\subsection}{\Needspace{5\baselineskip}}{}{\errmessage{cannot patch subsection}}

\tcbset{handoutbox/.style={enhanced, breakable, sharp corners, boxrule=0pt, left=7pt, right=7pt, top=5pt, bottom=5pt, before skip=8pt plus 2pt, after skip=8pt plus 2pt}}
\newtcolorbox{examplebox}{handoutbox, colback=hselight, colframe=hselight, before upper={\textcolor{hseblue}{\bfseries Example.}\ }}
\newtcolorbox{whybox}{handoutbox, colback=white, colframe=hseblue, leftrule=2.5pt, top=3pt, bottom=3pt, before upper={\textcolor{hseaccent}{\bfseries Why.}\ }}
\newtcolorbox{answerbox}{handoutbox, colback=hsewarm, colframe=hsewarm, before upper={\textcolor{hseaccent}{\bfseries Answer.}\ }}
\newtcolorbox{conventionbox}{handoutbox, colback=hselight, colframe=hseblue, leftrule=2.5pt, before upper={\textcolor{hseblue}{\bfseries Notation.}\par\smallskip}}

\newcommand{\handoutfigure}[3]{\par\medskip\noindent\begin{minipage}{\linewidth}\centering
  \includegraphics[width=#1\linewidth]{figures/#2.pdf}\par\smallskip
  \begin{minipage}{0.94\linewidth}\small\textcolor{hseblue}{\bfseries Figure.}\ #3\end{minipage}\end{minipage}\par\medskip}
\newcommand{\problemhead}[3]{\par\bigskip\noindent{\color{hseblue}\large\bfseries #1.\ #2}\par\nopagebreak
  \noindent{\small\itshape\color{hsegray}#3}\par\nopagebreak\smallskip}
\newtcolorbox{definitionbox}[1]{handoutbox, colback=white, colframe=hseblue, leftrule=2.5pt, before upper={\textcolor{hseblue}{\bfseries #1}\ }}
\newtcolorbox{theorembox}[1]{handoutbox, colback=hselight, colframe=hseblue, leftrule=2.5pt, before upper={\textcolor{hseblue}{\bfseries #1}\ }}
\newtcolorbox{proofbox}{handoutbox, colback=white, colframe=hsegray, leftrule=1.2pt, top=3pt, bottom=3pt, before upper={\textcolor{hsegray}{\bfseries\itshape Proof.}\ }}
\newtcolorbox{notebox}[1]{handoutbox, colback=hselight, colframe=hselight, before upper={\textcolor{hseblue}{\bfseries #1}\par\smallskip}}
\newtcolorbox{trapbox}{handoutbox, colback=hsewarm, colframe=hseaccent, leftrule=2.5pt, before upper={\textcolor{hseaccent}{\bfseries Trap.}\ }}
\newtcolorbox{selfcheckbox}{handoutbox, colback=white, colframe=hseaccent, leftrule=2.5pt, top=3pt, bottom=3pt, before upper={\textcolor{hseaccent}{\bfseries Check yourself.}\ }}
% a display wider than the line is set as a centred paragraph that may break after a relation
\newcommand{\handoutdisplay}[1]{\sbox0{$\displaystyle #1$}\ifdim\wd0>\linewidth
  \par{\centering\binoppenalty=10000\relpenalty=0\relax$\displaystyle #1$\par}\else\[#1\]\fi}
\newcommand{\selfcheckanswer}{\par\smallskip\textcolor{hsegray}{\bfseries Answer.}\ }

\setlist[itemize]{itemsep=3pt,topsep=4pt}
\setlist[enumerate]{itemsep=3pt,topsep=4pt}
\setlength{\parindent}{0pt}
\setlength{\parskip}{5pt plus 1pt}
\setlength{\emergencystretch}{2em}
\setcounter{tocdepth}{2}

\pagestyle{fancy}
\fancyhf{}
\fancyhead[L]{\small\color{hsegray}Optimization Methods}
\fancyhead[R]{\small\color{hsegray}Seminar 04}
\fancyfoot[L]{\small\color{hsegray}Evgeny Baulin}
\fancyfoot[C]{\small\thepage}
\fancyfoot[R]{\small\color{hsegray}MIT License}
\renewcommand{\headrulewidth}{0pt}
\renewcommand{\footrulewidth}{0pt}
\fancypagestyle{plain}{\pagestyle{fancy}}

\begin{document}

\begin{center}
  {\large\color{hsegray}Optimization Methods\par}
  \vspace{4pt}
  {\large\bfseries\color{hseaccent}Seminar 04\par}
  \vspace{8pt}
  {\LARGE\bfseries\color{hseblue}Accelerated Gradient Methods\par}
  \vspace{8pt}
  {\Large Theory handout\par}
  \vspace{12pt}
  HSE University, Faculty of Computer Science\par
  Academic year 2026–2027\par
  \vspace{8pt}
  Course lecturer: Andrey Ignatov\par
  Evgeny Baulin\par
  Telegram: \href{https://t.me/tarakan_tuc}{@tarakan\_tuc}\par
\end{center}

\noindent{\color{hseblue}\rule{\linewidth}{0.6pt}}

\tableofcontents

\clearpage


A slow gradient method may suffer from elongated geometry, expensive full gradients, or noisy estimates. This handout follows the lecture and computer seminar of Andrey Ignatov: momentum, coordinate descent, conjugate gradients, and stochastic least squares.

\begin{conventionbox}
$Q$ and $A$ are real symmetric positive definite matrices unless stated otherwise. $\mu$ and $L$ are their extreme eigenvalues, $\kappa=L/\mu$, $g=\nabla f$, $r=b-Ax$. For regression we use the \textbf{half} mean squared error, $F(w)=\|Xw-y\|^2/(2N)$; plain MSE has twice its gradient and Hessian.
\end{conventionbox}

% =====================================================================
\section{Geometry and a fair baseline}\label{sec:geometry}

\begin{definitionbox}{Definition.}
For $f(x)=\tfrac{1}{2}x^\top Qx-b^\top x$, the minimizer is $x^*=Q^{-1}b$. The error $e=x-x^*$ has objective gap $\tfrac{1}{2}e^\top Qe$ and gradient $Qe$. Write $Q=U\Lambda U^\top$ with $U^\top U=I$; all scalar recurrences below act on $z=U^\top e$.
\end{definitionbox}

\begin{theorembox}{Theorem.}
Fixed-step GD has $z_{k+1,i}=(1-\alpha\lambda_i)z_{k,i}$. Convergence from every start holds precisely when $0<\alpha<2/L$. The best constant quadratic step is $2/(L+\mu)$ with worst error factor $(\kappa-1)/(\kappa+1)$.
\end{theorembox}

\begin{proofbox}
Each coordinate is a geometric progression. Its magnitude tends to zero iff $|1-\alpha\lambda_i|<1$. To minimize the largest magnitude on $[\mu,L]$, balance the endpoint values $1-\alpha\mu=-(1-\alpha L)$.
\end{proofbox}

\begin{whybox}
With a large $\kappa$, the step needed to avoid exploding in the steep direction barely changes the flat direction. Orthogonal rotation preserves eigenvalues and condition number while moving the eigenvectors away from the coordinate axes.
\end{whybox}

\begin{trapbox}
An iteration is not a common unit of work. GD uses a full gradient, a coordinate step may need only one partial derivative, CG uses a matrix-vector product, and SGD uses a batch of sample gradients. Full-loss logging also reads all samples. Report these counts separately; wall time is a secondary, machine-dependent measurement.
\end{trapbox}

% =====================================================================
\section{Momentum on a quadratic}\label{sec:momentum}

\begin{definitionbox}{Definition.}
Heavy ball stores a displacement:
\handoutdisplay{
v_{k+1}=\beta v_k-\alpha g(x_k),\qquad x_{k+1}=x_k+v_{k+1}.
}
 Set $v_0=0$, equivalently $x_{-1}=x_0$, unless the task supplies a different previous point. Eliminating $v$ gives $x_{k+1}=x_k-\alpha g(x_k)+\beta(x_k-x_{k-1})$.
\end{definitionbox}

\begin{theorembox}{Theorem.}
For each eigenvalue $\lambda$, heavy ball has polynomial
\handoutdisplay{
p_\lambda(r)=r^2-(1+\beta-\alpha\lambda)r+\beta.
}
 For $0\le\beta<1$ the whole SPD quadratic is stable exactly when
\handoutdisplay{
0<\alpha<{2(1+\beta)\over L}.
}
 Stability here means convergence of the linear recurrence for every initial pair.
\end{theorembox}

\begin{proofbox}
For a real monic quadratic $r^2-ar+\beta$, both roots lie inside the unit circle iff $|\beta|<1$, $1-a+\beta>0$ and $1+a+\beta>0$. Substitute $a=1+\beta-\alpha\lambda$ to get $\alpha\lambda>0$ and $\alpha\lambda<2(1+\beta)$. The largest eigenvalue imposes the tightest restriction.
\end{proofbox}

\begin{theorembox}{Theorem.}
Polyak's parameters minimize the worst root radius over the spectral interval of an SPD quadratic:
\handoutdisplay{
\alpha_*={4\over(\sqrt L+\sqrt\mu)^2},\quad \beta_*=\left({\sqrt L-\sqrt\mu\over\sqrt L+\sqrt\mu}\right)^2,\quad \rho_*={\sqrt\kappa-1\over\sqrt\kappa+1}.
}
 The resulting scale of iteration complexity is $O(\sqrt\kappa\log(1/\varepsilon))$, compared with $O(\kappa\log(1/\varepsilon))$ for GD.
\end{theorembox}

\begin{proofbox}
Put $s=\sqrt{\beta_*}=\rho_*$. The endpoint coefficients satisfy $1+\beta_*-\alpha_*\mu=2s$ and $1+\beta_*-\alpha_*L=-2s$. At interior eigenvalues the discriminant is nonpositive, so the complex-conjugate roots have modulus $s$. At the endpoints they are repeated. Repeated-root modes have the form $(c+dk)(\pm s)^k$; a pure bound $\|e_k\|\le s^k\|e_0\|$ does not follow.
\end{proofbox}

\begin{trapbox}
These optimal parameters are a quadratic result. Do not transfer the same global guarantee to every smooth strongly convex nonlinear objective. Nor does a smaller root radius require the objective to decrease at each early step.
\end{trapbox}

\problemhead{A1}{Heavy ball: parameters and the first steps}{On the board}
\textbf{Statement.} For $f(x)=\tfrac{1}{2}x^\top Qx$, $Q=\operatorname{diag}(1,\ 9)$, use $x_{-1}=x_0=(1,\ 1)$ and $x_{k+1}=x_k-\alpha Qx_k+\beta(x_k-x_{k-1})$. Find Polyak's optimal quadratic parameters and the spectral radius. Calculate $x_1,x_2$ and their objective values. Must the objective fall at every step?

\textbf{Solution.}
\begingroup\sloppy
\begin{enumerate}[label=\textbf{\arabic*.}, leftmargin=*]
  \item Here $\mu=1$, $L=9$, $\kappa=9$. For an SPD quadratic,
\handoutdisplay{
\alpha_*={4\over(\sqrt L+\sqrt\mu)^2}=0.25,\qquad \beta_*=\left({\sqrt L-\sqrt\mu\over\sqrt L+\sqrt\mu}\right)^2=0.25.
}
 The maximum root modulus is $\rho=\sqrt{\beta_*}=0.5$.
  \item Along an eigenvector with eigenvalue $\lambda$,
\handoutdisplay{
z_{k+1}=(1+\beta-\alpha\lambda)z_k-\beta z_{k-1}.
}
 The characteristic polynomial is $r^2-(1+\beta-\alpha\lambda)r+\beta$. At the two endpoints its roots are repeated: $+\rho$ and $-\rho$.
  \item The first displacement has no momentum because $x_0-x_{-1}=0$. Thus $x_1=(0.75,\ -1.25)$; the next step includes $\beta(x_1-x_0)$ and gives $x_2=(0.5,\ 1)$.
  \item $f(x_0)=5$, $f(x_1)=7.3125$, $f(x_2)=4.625$. The objective rises at the first step, although every characteristic root is inside the unit circle. Root radius describes eventual decay; repeated roots can add a factor proportional to $k$.
\end{enumerate}
\endgroup
\begin{answerbox}
$\alpha_*=0.25$, $\beta_*=0.25$, $\rho=0.5$. Stable does not mean monotone.
\end{answerbox}

\problemhead{A2}{The momentum stability boundary}{On your own}
\textbf{Statement.} On $f(z)=\tfrac{1}{2}\lambda z^2$, take $\lambda=16$ and $\beta=0.5$. Derive the strict interval of stable positive steps. Report its upper endpoint and the largest root modulus for $\alpha=0.2$. Is that trial stable? Treat equality at the endpoint separately.

\textbf{Solution.}
\begingroup\sloppy
\begin{enumerate}[label=\textbf{\arabic*.}, leftmargin=*]
  \item For $a=1+\beta-\alpha\lambda$, the root test reduces to $|\beta|<1$, $\alpha\lambda>0$, and $2(1+\beta)-\alpha\lambda>0$. With the stated nonnegative momentum, $0<\alpha<2(1+\beta)/\lambda=0.1875$.
  \item At the endpoint the polynomial factors as $(r+1)(r+\beta)$, so one root is $-1$. A generic starting pair retains an oscillating component. At the proposed larger step the largest modulus is $1.3217>1$, so some starts diverge.
\end{enumerate}
\endgroup
\begin{answerbox}
Upper endpoint: $0.1875$; trial radius: $1.3217$. The stable interval is open.
\end{answerbox}

\handoutfigure{0.95}{fig_momentum}{A nontrivial rotation makes the steep and flat directions visible without aligning them to the axes.}

\handoutfigure{0.9}{fig_stability}{The analytic boundary is strict; a finite run also depends on the reduction threshold.}

% =====================================================================
\section{Look-ahead and restart}\label{sec:lookahead}

\begin{definitionbox}{Definition.}
The constant-momentum look-ahead form is
\handoutdisplay{
y_k=x_k+\beta(x_k-x_{k-1}),\qquad x_{k+1}=y_k-\alpha g(y_k).
}
 On a scalar quadratic it gives
\handoutdisplay{
z_{k+1}=(1-\alpha\lambda)\big((1+\beta)z_k-\beta z_{k-1}\big).
}
 Its polynomial differs from heavy ball: $r^2-(1+\beta)(1-\alpha\lambda)r+\beta(1-\alpha\lambda)$.
\end{definitionbox}

\begin{notebox}{Two acceleration settings}
For a known strongly convex smooth objective, a standard accelerated choice is $\alpha=1/L$, $\beta=(\sqrt L-\sqrt\mu)/(\sqrt L+\sqrt\mu)$. In the merely convex case, acceleration uses a changing extrapolation coefficient; the $O(1/k^2)$ objective bound is not a claim about arbitrary constant momentum.
\end{notebox}

\begin{whybox}
If the extrapolated step points the wrong way, discard its momentum and retry a gradient step from the current iterate. An objective restart tests whether the proposed objective exceeds the current one. Label the restarted method separately and count the additional gradient evaluation.
\end{whybox}

\problemhead{A3}{Where is the gradient evaluated?}{On your own}
\textbf{Statement.} For $f(x)=\tfrac{1}{2}\lambda x^2$, use $\lambda=4$, $x_k=2$, $x_{k-1}=3$, $\alpha=0.1$, $\beta=0.5$. Compute the next heavy-ball iterate and the next look-ahead iterate: $y_k=x_k+\beta(x_k-x_{k-1})$, $x_{k+1}=y_k-\alpha\nabla f(y_k)$. Explain why they differ.

\textbf{Solution.}
\begingroup\sloppy
\begin{enumerate}[label=\textbf{\arabic*.}, leftmargin=*]
  \item The extrapolated point is $y_k=1.5$. Heavy ball uses $\nabla f(x_k)=8$ and gives $x_k+\beta(x_k-x_{k-1})-\alpha\nabla f(x_k)=0.7$.
  \item The look-ahead update takes $\nabla f(y_k)=6$, giving $y_k-\alpha\nabla f(y_k)=0.9$. Using a gradient at $x_k$ in both routines silently makes them identical.
\end{enumerate}
\endgroup
\begin{answerbox}
Heavy ball: $0.7$. Look-ahead: $0.9$.
\end{answerbox}

% =====================================================================
\section{Coordinate descent}\label{sec:coordinates}

\begin{definitionbox}{Definition.}
A coordinate Lipschitz constant $L_i>0$ gives
\handoutdisplay{
f(x+te_i)\le f(x)+tg_i+\tfrac{1}{2}L_it^2.
}
 Choosing $t=-g_i/L_i$ yields
\handoutdisplay{
f(x^+)\le f(x)-{g_i^2\over2L_i}.
}
 For a quadratic $L_i=Q_{ii}$ and the inequality is an equality.
\end{definitionbox}

\begin{proofbox}
For a quadratic, expand the slice exactly: $f(x+te_i)=f(x)+tg_i+Q_{ii}t^2/2$. Differentiating in $t$ gives its unique minimizer. Substitution gives the decrease. After modifying $x_i$, the next derivative must use the modified vector.
\end{proofbox}

\begin{notebox}{Three selectors}
Cyclic selection visits each coordinate in order. Uniform random selection samples one coordinate per update. Greedy selection maximizes $|g_i|/\sqrt{L_i}$, since its square is twice the guaranteed decrease. In $d$ dimensions, $d$ coordinate updates are one equivalent pass, but a pass is not necessarily as cheap as one full gradient.
\end{notebox}

\begin{theorembox}{Theorem.}
If $f$ is $\mu$-strongly convex and all coordinates have $L_i\le L_{\max}$, uniform coordinate selection gives
\handoutdisplay{
\mathbb E[f(x^+)-f^*\mid x]\le\left(1-{\mu\over dL_{\max}}\right)(f(x)-f^*).
}
 This is a bound per coordinate update, under fresh independent coordinate sampling.
\end{theorembox}

\begin{proofbox}
Average $g_i^2/(2L_i)$ over the coordinates: the expected decrease is at least $\|g\|^2/(2dL_{\max})$. Strong convexity implies $\|g\|^2\ge2\mu(f-f^*)$; substitute into the decrease estimate.
\end{proofbox}

\problemhead{B1}{One coordinate at a time}{On the board}
\textbf{Statement.} For $f(x)=\tfrac{1}{2}x^\top Qx$, $Q=\begin{pmatrix} 5 & -4 \\ -4 & 5 \end{pmatrix}$, start at $x_0=(1,\ 1)$. Minimize first over coordinate $1$, then over coordinate $2$, keeping the other fixed. Derive the exact coordinate step and calculate the objective decrease after each update.

\textbf{Solution.}
\begingroup\sloppy
\begin{enumerate}[label=\textbf{\arabic*.}, leftmargin=*]
  \item Along coordinate $i$, $f(x+te_i)=f(x)+tg_i+\tfrac{1}{2}Q_{ii}t^2$. Its minimizer is $t=-g_i/Q_{ii}$, and the decrease is $g_i^2/(2Q_{ii})$. Here $L_i=Q_{ii}$, even though the eigenvalues are $(1,\ 9)$.
  \item $g(x_0)=(1,\ 1)$. After coordinate $1$: $x=(0.8,\ 1)$, $f=0.9$, decrease $0.1$. Recompute the gradient at this point before choosing the next step.
  \item After coordinate $2$: $x=(0.8,\ 0.64)$, $f=0.576$, decrease $0.324$. These are two coordinate updates, or one full pass in this two-dimensional problem.
\end{enumerate}
\endgroup
\begin{answerbox}
$x_1=(0.8,\ 1)$, $x_2=(0.8,\ 0.64)$. Each update minimizes the current one-dimensional slice.
\end{answerbox}

\problemhead{B2}{Choose the coordinate by its decrease}{On your own}
\textbf{Statement.} For $Q=\operatorname{diag}(1,\ 9)$ and $x=(2,\ 1)$, compute both gradient entries. Which coordinate maximizes the exact decrease $g_i^2/(2Q_{ii})$? Report the coordinate number, the decrease, and the new objective value. Explain the score $|g_i|/\sqrt{Q_{ii}}$.

\textbf{Solution.}
\begingroup\sloppy
\begin{enumerate}[label=\textbf{\arabic*.}, leftmargin=*]
  \item The gradient is $(2,\ 9)$. The normalized scores are $(2,\ 3)$, so coordinate $2$ gives the largest exact improvement.
  \item Updating this coordinate reduces $f$ by $4.5$ and leaves $f=2$. Maximizing $|g_i|/\sqrt{Q_{ii}}$ is equivalent to maximizing its square and therefore the exact coordinate decrease.
\end{enumerate}
\endgroup
\begin{answerbox}
Coordinate $2$; decrease $4.5$; new objective $2$.
\end{answerbox}

% =====================================================================
\section{Conjugate gradients}\label{sec:cg}

\begin{definitionbox}{Definition.}
For $A\succ0$, minimizing $f(x)=\tfrac{1}{2}x^\top Ax-b^\top x$ is equivalent to solving $Ax=b$. Directions $p_i,p_j$ are $A$-conjugate if $p_i^\top Ap_j=0$. The residual $r=b-Ax$ is the negative gradient.
\end{definitionbox}

\begin{definitionbox}{Definition.}
Initialize $r_0=b-Ax_0$, $p_0=r_0$. At each step compute
\handoutdisplay{
\begin{gathered}
q_k=Ap_k,\quad \alpha_k={r_k^\top r_k\over p_k^\top q_k}, \\[6pt]
x_{k+1}=x_k+\alpha_kp_k,\quad r_{k+1}=r_k-\alpha_kq_k, \\[6pt]
\beta_k={r_{k+1}^\top r_{k+1}\over r_k^\top r_k},\quad p_{k+1}=r_{k+1}+\beta_kp_k.
\end{gathered}
}
 Check convergence before forming the next coefficient if the new residual is zero.
\end{definitionbox}

\begin{theorembox}{Theorem.}
In exact arithmetic, distinct residuals are orthogonal and search directions are $A$-conjugate. CG minimizes the quadratic over $x_0+\mathcal K_k(A,r_0)$, where $\mathcal K_k=\operatorname{span}\{r_0,Ar_0,\ldots,A^{k-1}r_0\}$. It terminates in at most $d$ steps in dimension $d$.
\end{theorembox}

\begin{whybox}
Minimizing along a new conjugate direction does not undo the minimizations along earlier directions: $p_i^\top A(x+tp_j)=p_i^\top Ax$ for $i\ne j$. Once the search directions span the full space, the gradient is orthogonal to the entire space and must vanish.
\end{whybox}

\begin{theorembox}{Theorem.}
With the energy norm $\|e\|_A=\sqrt{e^\top Ae}$,
\handoutdisplay{
\|e_k\|_A\le2\left({\sqrt\kappa-1\over\sqrt\kappa+1}\right)^k\|e_0\|_A.
}
 Since $f(x_k)-f^*=\|e_k\|_A^2/2$, the corresponding objective ratio has the squared factor. Eigenvalue clustering can make the observed convergence much faster than this worst-case interval bound.
\end{theorembox}

\begin{trapbox}
Use $\|r\|/\|b\|$ when $b\ne0$. For $b=0$ use a documented absolute threshold. Check symmetry and positive definiteness for small dense examples; guard against nonpositive curvature, nonfinite values and a maximum iteration limit. The true final residual is $b-Ax$, and its verification costs an additional matrix-vector product.
\end{trapbox}

\problemhead{B3}{Two conjugate-gradient steps}{On your own}
\textbf{Statement.} Solve $Ax=b$ by CG with $A=\begin{pmatrix} 4 & 1 \\ 1 & 3 \end{pmatrix}$, $b=(1,\ 2)$, $x_0=(0,\ 0)$. Start with $r_0=b-Ax_0$, $p_0=r_0$. Compute $\alpha_0$, $\beta_0$, and $x_2$. Check $r_0^\top r_1=0$ and $p_0^\top Ap_1=0$. Compare with the direct solution.

\textbf{Solution.}
\begingroup\sloppy
\begin{enumerate}[label=\textbf{\arabic*.}, leftmargin=*]
  \item The first line minimization gives $\alpha_0=0.25$, $x_1=(0.25,\ 0.5)$, $r_1=(-0.5,\ 0.25)$. Its inner product with $r_0$ is $0$.
  \item $\beta_0=0.0625$, $p_1=r_1+\beta_0p_0=(-0.4375,\ 0.375)$. The directions satisfy $p_0^\top Ap_1=0$. The next step is $\alpha_1=0.3636$.
  \item $x_2=x^*=(0.0909,\ 0.6364)$ and $r_2=0$ in exact arithmetic. For an SPD matrix CG terminates in at most the dimension, but rounding may require more steps. The repeated matrix-vector product $Ap_k$ must be reused in the denominator and residual update.
\end{enumerate}
\endgroup
\begin{answerbox}
$\alpha_0=0.25$, $\beta_0=0.0625$, $x_2=(0.0909,\ 0.6364)$.
\end{answerbox}

\handoutfigure{0.95}{fig_directions}{GD, cyclic coordinate descent and CG on the same SPD quadratic, with their different work units.}

% =====================================================================
\section{Mini-batches and the noise floor}\label{sec:sgd}

\begin{definitionbox}{Definition.}
For $F(w)=N^{-1}\sum_i\ell_i(w)$, a fresh uniform index gives $\mathbb E[\nabla\ell_i(w)\mid w]=\nabla F(w)$. A batch gradient is the mean of its sample gradients. If $B$ samples are independent and each has variance $\sigma^2$, the batch variance is $\sigma^2/B$.
\end{definitionbox}

\begin{proofbox}
The sample mean is linear, so its expectation is the full mean. Expand the squared centred average: independent cross terms have zero expectation, leaving $B$ diagonal terms divided by $B^2$. Without replacement, covariance terms are negative and yield the finite-population correction $(N-B)/(N-1)$.
\end{proofbox}

\begin{notebox}{Sampling convention}
Shuffling without replacement each epoch is a useful implementation choice. However, an iterate depends on the earlier batches, so conditional unbiasedness of every later batch is not automatic. State whether a theoretical argument assumes fresh independent batches or random reshuffling.
\end{notebox}

\begin{theorembox}{Theorem.}
For an $L$-smooth objective and an unbiased gradient with conditional noise variance at most $\sigma^2/B$, the descent lemma gives
\handoutdisplay{
\mathbb E[F(w^+)\mid w]\le F(w)-\alpha(1-L\alpha/2)\|\nabla F(w)\|^2+{L\alpha^2\sigma^2\over2B}.
}
 The positive noise term explains why a constant step can leave a floor.
\end{theorembox}

\begin{proofbox}
Insert $w^+=w-\alpha g$ into the smoothness upper model. Use $\mathbb E[g\mid w]=\nabla F(w)$ and $\mathbb E[\|g\|^2\mid w]\le\|\nabla F(w)\|^2+\sigma^2/B$. Collect the gradient and variance terms.
\end{proofbox}

\problemhead{C1}{A batch estimates a gradient}{On the board}
\textbf{Statement.} Let $\ell_i(w)=\tfrac{1}{2}(w-a_i)^2$, $a=(0,\ 2,\ 4,\ 6)$, and $F(w)$ be the mean loss. At $w_0=0$, compare a full GD step with a batch containing targets $(0,\ 2)$, using $\alpha=0.25$. Compute the variance of one uniformly sampled gradient. How does averaging two samples change it with and without replacement?

\textbf{Solution.}
\begingroup\sloppy
\begin{enumerate}[label=\textbf{\arabic*.}, leftmargin=*]
  \item At the start, the full gradient is $-3$ and the selected batch gradient is $-1$. Hence $w_{\mathrm{full}}=0.75$ and $w_B=0.25$. The batch is an estimate, rather than the exact full gradient.
  \item The minimizer is the mean target, $w^*=3$. Subtract the full gradient from each sample gradient and square; their mean is $\sigma^2=5$. Averaging two independent draws gives $2.5$.
  \item For a uniformly chosen subset of size $B$ without replacement, $\operatorname{Var}(g_B)=\sigma^2(N-B)/(B(N-1))$, so this batch has variance $1.6667$. Inside a shuffled epoch, later batches are conditioned on earlier draws; do not silently reuse conditional independence.
\end{enumerate}
\endgroup
\begin{answerbox}
$g_B=-1$, $w_B=0.25$, $\sigma^2=5$.
\end{answerbox}

\problemhead{C2}{The constant-step noise floor}{On your own}
\textbf{Statement.} For $F(w)=\tfrac{1}{2}\mu w^2$, the batch gradient is $\mu w+\xi_k$, where noises are independent, centred and have variance $\sigma^2/B$. Take $\mu=2$, $\sigma^2=9$, $B=3$, $\alpha=0.1$. Derive the recurrence for $\mathbb E w_k^2$ and its stationary value. Report stationary variance and expected objective gap.

\textbf{Solution.}
\begingroup\sloppy
\begin{enumerate}[label=\textbf{\arabic*.}, leftmargin=*]
  \item $w_{k+1}=(1-\alpha\mu)w_k-\alpha\xi_k$. Independence and zero mean give
\handoutdisplay{
\mathbb E w_{k+1}^2=(1-\alpha\mu)^2\mathbb E w_k^2+\alpha^2\sigma^2/B.
}
 Here the mean contracts by $0.8$.
  \item For $0<\alpha\mu<2$, solve the fixed-point equation:
\handoutdisplay{
V_\infty={\alpha\sigma^2\over B\mu(2-\alpha\mu)}=0.0833.
}
 The limiting mean is zero, so this second moment is a variance; $\mathbb E[F(w)-F^*]=\mu V_\infty/2=0.0833$.
  \item A smaller step or a larger batch lowers the floor. Its RMS radius is $0.2887$, not the variance itself. A shrinking learning rate can remove this constant floor under suitable assumptions.
\end{enumerate}
\endgroup
\begin{answerbox}
$V_\infty=0.0833$, expected gap $0.0833$.
\end{answerbox}

% =====================================================================
\section{Schedules and reproducible experiments}\label{sec:schedules}

\begin{definitionbox}{Definition.}
Use the \textbf{update count} $k$, starting from zero, to define the schedules: constant $\alpha_k=\alpha_0$, inverse square root $\alpha_k=\alpha_0/\sqrt{1+k}$, inverse $\alpha_k=\alpha_0/(1+k)$. An epoch counter defines a different schedule; do not mix the two conventions.
\end{definitionbox}

\begin{notebox}{What a schedule guarantees}
Under the usual stochastic-approximation assumptions, $\sum_k\alpha_k=\infty$ and $\sum_k\alpha_k^2<\infty$ are sufficient step conditions. Inverse steps meet both; inverse-square-root steps fail square summability. That does not make the latter unusable for finite-budget convex experiments, but it rules out quoting the same theorem without further work.
\end{notebox}

\begin{trapbox}
An inverse schedule can decay too quickly before the method traverses an elongated valley. Compare the same initial rate, update convention, data, seeds and work budget. A schedule that wins after a short run may lose after a longer one.
\end{trapbox}

\begin{notebox}{Epoch protocol}
At each epoch draw a permutation, take slices of size $B$, and keep the last slice even if it is shorter. Use its actual length in the mean gradient. Log the complete objective after finishing the epoch. Repeat across $5$ random seeds and show the median and interquartile interval.
\end{notebox}

\handoutfigure{0.9}{fig_sgd}{Use the same sample-gradient budget for all schedules; the shaded bands show seed variability.}

% =====================================================================
\section{Feature standardization}\label{sec:scaling}

\begin{definitionbox}{Definition.}
For centred $X$, let $S=\operatorname{diag}(s_i)$ with positive training standard deviations and define $Z=XS^{-1}$. Its least-squares Hessian is $H_Z=S^{-1}HS^{-1}$ with $H=X^\top X/N$. Prediction equality $Zu=Xw$ requires $w=S^{-1}u$.
\end{definitionbox}

\begin{whybox}
Columns with different units stretch the quadratic. Scaling can remove this part of the condition number. It does not remove near-linear dependence between columns, so the standardized Hessian can still be ill-conditioned. Constant columns need a separate convention rather than division by zero.
\end{whybox}

\begin{notebox}{Restore original units}
With means $m$, standardized prediction $c+(x-m)^\top S^{-1}u$ becomes $b+x^\top w$, where $w=S^{-1}u$ and $b=c-m^\top w$. Estimate $m,S$ on training data only. A grid search for the largest usable SGD step is an empirical finite-run test, distinct from the full-gradient boundary $2/\lambda_{\max}(H)$.
\end{notebox}

\problemhead{C3}{Scale features, then restore units}{On your own}
\textbf{Statement.} For centred least squares $F(w)=\|Xw-y\|^2/(2N)$, let $H=X^\top X/N=\begin{pmatrix} 1 & 0 \\ 0 & 100 \end{pmatrix}$. Standardize with $Z=XS^{-1}$, $S=\operatorname{diag}(1,\ 10)$. Find the condition numbers before and after scaling and the full-GD stability boundaries. If the standardized fit returns $u=(2,\ 30)$, what are the coefficients in the original units?

\textbf{Solution.}
\begingroup\sloppy
\begin{enumerate}[label=\textbf{\arabic*.}, leftmargin=*]
  \item $H_Z=S^{-1}HS^{-1}=I$. Thus $\kappa(H)=100$, $\kappa(H_Z)=1$. The upper full-GD step endpoints are $2/L=0.02$ and $2$, respectively; both boundaries are excluded.
  \item Convert the coefficients back: $w=S^{-1}u=(2,\ 3)$. For non-centred features with mean $m$, a standardized intercept $c$ becomes $c-m^\top w$. Equal coefficients in the two coordinate systems are not expected.
  \item Scaling fixes unequal units in this diagonal example. Strong correlations can remain after standardization, and a full-Hessian bound is not an exact stability guarantee for every individual stochastic batch.
\end{enumerate}
\endgroup
\begin{answerbox}
$\kappa$: $100\to1$; $w=(2,\ 3)$.
\end{answerbox}

\handoutfigure{0.9}{fig_scaling}{Original and standardized least-squares losses with identical SGD settings.}

% =====================================================================
\section{Implementation checklist and exercises}\label{sec:implementation}

\begin{itemize}
  \item Return an independent copy of every saved iterate; appending the same mutable array erases the path.
  \item Stop on the current gradient or residual, guard nonfinite values, and document the maximum iteration limit.
  \item Report initial setup, optimization and diagnostic work separately. A greedy coordinate scan computes more than the selected derivative.
  \item Compare NumPy implementations against exact small examples and direct linear algebra, then run the larger experiments.
\end{itemize}

\begin{selfcheckbox}
Why can an analytically stable pair fail a finite-run speed test?
\selfcheckanswer A root modulus close to one and an initial transient may keep the objective above the chosen finite-run threshold.
\end{selfcheckbox}

\begin{selfcheckbox}
Why is one CG iteration not directly comparable to one coordinate update?
\selfcheckanswer CG computes a matrix-vector product, while a coordinate update can use one partial derivative. The real cost depends on sparsity and on the implementation of coordinate selection.
\end{selfcheckbox}

\end{document}
